% =====================================================================
%  CAPITULO I - REVISION CRITICA DE LA LITERATURA
%  Cada estudio explica lo que aporta al proyecto; el aporte del proyecto va al final.
%  El Aporte del proyecto (A1..A7) se enuncia al final, en la Seccion 1.5.
%  Aportes numerados A1..A7, secuenciales en orden de lectura.
%  Todas las cifras citadas provienen de las tablas de los articulos
%  originales, con la procedencia indicada en nota al pie.
% =====================================================================

\chapter{REVISIÓN CRÍTICA DE LA LITERATURA}
\label{cap:revision}

Las cuatro publicaciones revisadas se corresponden con los objetivos específicos
de la tesis (Tabla~\ref{tab:obj-vs-revision}). De cada una se explica primero lo
que \textbf{aporta al proyecto}: qué hizo, qué resultado obtuvo y por qué sirve.
Al final, en la Sección~\ref{sec:aporte}, se enuncian los siete \textbf{aportes
del proyecto}, es decir, lo que esta tesis añade a partir de esos estudios.

\begin{table}[H]
\centering
\setstretch{1}\footnotesize
\renewcommand{\arraystretch}{1.4}
\caption[Aportes del proyecto]{Inventario de los siete aportes del proyecto. Cada uno nace de un estudio
revisado y termina en un resultado propio que se puede verificar.}
\label{tab:obj-vs-revision}
\begin{tabularx}{\linewidth}{@{}l>{\raggedright\arraybackslash\hsize=0.90\hsize}X>{\raggedright\arraybackslash\hsize=0.83\hsize}X>{\raggedright\arraybackslash\hsize=1.27\hsize}X@{}}
\hline
 & \textbf{Aporte} & \textbf{Nace de} & \textbf{Resultado verificable} \\
\hline
\multicolumn{4}{@{}l}{\textbf{OE1. Integración de Sonata en el módulo de prototipos}}\\
A1 & Sustitución de $f_{\text{geo}}$ por $f_{\mathrm{ssl}}$
   & HyperNCD (\S\ref{sec:hyperncd}), Sonata (\S\ref{sec:sonata})
   & Ramas congeladas y fusión posterior con métricas comparables \\
A2 & Auditoría del código y corrección de la línea base
   & HyperNCD (\S\ref{sec:hyperncd}) y su repositorio
   & Tabla artículo frente a implementación y tres decisiones metodológicas \\
A4 & Puente entre serialización y rejilla dispersa
   & PTv3 (\S\ref{sec:ptv3})
   & \emph{Cache} por sección y cota de aceleración propia \\
\hline
\multicolumn{4}{@{}l}{\textbf{OE2. Calibración del híbrido}}\\
A5 & Normalización de escala entre ramas heterogéneas
   & BRDL (\S\ref{sec:brdl})
   & Distribución de $S_b$ antes y después \\
A6 & Calibración conjunta de $(\alpha,M,\epsilon)$ y no-aditividad
   & BRDL (\S\ref{sec:brdl})
   & Grilla conjunta, curvas de sensibilidad, configuración fijada \\
A7 & Robustez de la configuración entre secciones
   & BRDL (\S\ref{sec:brdl})
   & Métrica por sección sin reajuste \\
\hline
\multicolumn{4}{@{}l}{\textbf{OE3. Cuantificación de la mejora y rango de operación}}\\
A3 & Brecha cerrada $G$ y atribución por etapas
   & Sonata (\S\ref{sec:sonata})
   & $G$ formulada, medición intermedia, PCA \\
\hline
\end{tabularx}
\end{table}

% =====================================================================
\section[HyperNCD]{Zhang, Wu, Li, Han y Shen (2026): HyperNCD, razonamiento por
hipergrafo}
\label{sec:hyperncd}

Presentado en CVPR 2026~\cite{zhang2026hyperncd}, es el método sobre el que se
plantea \textbf{OE1}. Representa las relaciones entre clases con un
\textbf{hipergrafo}. En un grafo común una arista une dos nodos; en un hipergrafo una hiperarista puede
unir varios a la vez, lo que permite expresar que varias clases comparten
una misma estructura geométrica. Los nodos son los prototipos de clase,
construidos sobre un descriptor compuesto: uno semántico $f_{\text{sem}}$, que
produce la red supervisada, y uno geométrico $f_{\text{geo}}$, calculado con una
fórmula fija.

Sobre SemanticKITTI obtiene entre 11,8 y 47,5 de mIoU en clases novedosas según
la partición, frente a 49,8 del equivalente supervisado
\cite[Tabla~1]{zhang2026hyperncd}. Más importante es su estudio de ablación, que apaga una pieza a la vez. En la
partición~0 de SemanticPOSS, desde una red base de 23,2 de mIoU, el prototipo
geométrico aporta $+4{,}2$ puntos y el hipergrafo $+9{,}6$
(Tabla~\ref{tab:ablacion-hyperncd}).

\begin{table}[!htbp]
\centering
\setstretch{1}\small
\caption[Ablación de HyperNCD en SemanticPOSS]{Ablación de HyperNCD sobre la partición~0 de SemanticPOSS, según la
Tabla~3 de \cite{zhang2026hyperncd}. Sostiene cuantitativamente la hipótesis de esta tesis.}
\label{tab:ablacion-hyperncd}
\begin{tabularx}{\linewidth}{@{}>{\raggedright\arraybackslash\hsize=1.00\hsize}Xcc@{}}
\hline
Configuración & mIoU promedio & Ganancia sobre la base \\
\hline
Red base (MinkowskiUNet-34C modificada) & 23,2 & referencia \\
Solo prototipo consciente de la geometría & 27,4 & $+4{,}2$ \\
Solo hipergrafo & 32,8 & $+9{,}6$ \\
Prototipo geométrico e hipergrafo & 34,2 & $+11{,}0$ \\
Método completo, con hiperaristas dinámicas & 36,5 & $+13{,}3$ \\
\hline
\end{tabularx}
\end{table}

La ablación muestra además que el descriptor geométrico aporta menos de la mitad
que el hipergrafo, una diferencia que los autores no discuten y que sugiere que
una representación aprendida tendría margen para aportar más.

Su límite de fondo es ese descriptor: se calcula de forma analítica con los 15
vecinos más cercanos \cite{zhang2026hyperncd}, es una fórmula fija y no se adapta
al dominio.

% =====================================================================
\section[Sonata]{Wu et al. (2025): Sonata y el atajo geométrico}
\label{sec:sonata}

Presentado en CVPR 2025~\cite{wu2025sonata}, aporta el codificador
auto-supervisado que \textbf{OE1} integra y, por su metodología de medición,
sirve a \textbf{OE3}. Identifica y ataca el \textbf{atajo geométrico}: cuando un
modelo se entrena sin etiquetas sobre nubes de puntos, la representación tiende a colapsar
sobre pistas geométricas de bajo nivel fáciles de leer, como la normal de la
superficie o la altura del punto, en lugar de aprender semántica. Sonata lo
ataca oscureciendo la información espacial dentro de un esquema de
auto-destilación entrenado sobre 140 mil escenas.

Su evidencia se mide por \emph{linear probing}: se congelan todos los pesos del
codificador y se entrena encima solo un clasificador lineal. Sobre ScanNet el
desempeño pasa de 5,6 a 72,5 de mIoU (Tabla~\ref{tab:linprobe}).

\begin{table}[!htbp]
\centering
\setstretch{1}\small
\caption[\emph{Linear probing} sobre ScanNet]{\emph{Linear probing} sobre ScanNet (validación), según la
Tabla~5 de \cite{wu2025sonata}. Codificador congelado y una sola capa lineal.}
\label{tab:linprobe}
\begin{tabular}{lc}
\hline
Método & mIoU (\%) \\
\hline
PointContrast & 5,6 \\
CSC & 12,6 \\
MSC & 21,8 \\
DINOv2, agregado de 2D a 3D & 63,1 \\
\textbf{Sonata} & \textbf{72,5} \\
\hline
\end{tabular}
\end{table}

Esto sirve por dos razones. La primera es el resultado: la señal geométrica de
bajo nivel es un descriptor débil en igualdad de condiciones, y existe un
codificador público, entrenado y congelable, que la reemplaza sin
preentrenamiento propio. Es la pieza que ocupa el lugar de $f_{\text{geo}}$ en el
aporte A1. La segunda es de método: el \emph{linear probing} aísla la
representación, y ese principio permite separar después lo que aporta ella de lo
que aporta el razonamiento.

El trabajo declara tres limitaciones que acotan el riesgo del proyecto:

\begin{enumerate}
\item \textit{Dominio de preentrenamiento.} El modelo se preentrena sobre escenas
interiores densas, entre ellas ScanNet, S3DIS y Structured3D
\cite{wu2025sonata}. Una catedral cae del lado favorable de esa frontera, pero el
cambio de dominio sigue siendo el principal riesgo declarado.

\item \textit{Régimen congelado en exteriores.} En escaneo láser urbano el
\emph{linear probing} queda por debajo de una PTv3 supervisada, aunque con ajuste
fino completo la supera \cite[Tabla~8]{wu2025sonata}. La desventaja es del
régimen congelado, que es el que adopta este proyecto, y no de la representación.

\item \textit{Vocabularios grandes.} Con 200 clases cae a 29,3 de mIoU en
ScanNet200, frente a 72,5 con 20 en ScanNet, lo que fija el límite superior de
clases novedosas simultáneas del rango favorable.
\end{enumerate}

% =====================================================================
\section[Point Transformer V3]{Wu et al. (2024): Point Transformer V3 y el
problema de integración}
\label{sec:ptv3}

Presentado en CVPR 2024~\cite{wu2024ptv3}, es la red sobre la que está construido
Sonata: el codificador que \textbf{OE1} emplea congelado es una PTv3 con pesos
auto-supervisados. Sustituye la vecindad calculada por vecinos más cercanos por
una vecindad \emph{serializada}: los puntos se ordenan siguiendo una curva de
llenado de espacio, la secuencia resultante se corta en parches contiguos y la
atención se calcula dentro de cada parche.

Las dos ramas del híbrido parten de las mismas coordenadas pero definen la
vecindad de maneras distintas. HyperNCD opera sobre una \textbf{rejilla dispersa
de Minkowski}, donde el vecino de un punto es el que cae dentro del núcleo de la
convolución, o sea, adyacencia métrica en el volumen. Sonata, a través de PTv3,
opera sobre una \textbf{secuencia unidimensional}, donde el vecino es el que
quedó al lado dentro de un parche de 1024 elementos de la curva. Una curva de
llenado de espacio preserva la localidad de forma aproximada y no exacta: dos
puntos contiguos en el volumen pueden quedar separados en la secuencia si caen a
ambos lados de un giro. Pasar la salida de una rama a la otra exige absorber ese
desajuste de forma explícita.

Se suma una segunda incompatibilidad, de ingeniería y no de arquitectura. La
inspección del código de ambos grupos muestra que las dos ramas exigen
\textbf{pilas de dependencias mutuamente excluyentes}: MinkowskiEngine contra
CUDA 11.3 y PTv3 contra CUDA 12.4. Integrar las dos implementaciones en un único
proceso no es posible sin renunciar a una, de modo que el punto de fusión no puede
ser un objeto en memoria compartido.

El costo del híbrido no viene de añadir un transformador, sino de ejecutarlo
dentro del bucle de entrenamiento. Lo que lo hace resoluble es que el codificador
se emplea congelado y solo en inferencia: su salida no depende de la época ni del
optimizador, así que se calcula una sola vez por sección.

% =====================================================================
\section[BRDL]{Su, Chen y Wang (2025): BRDL, calibración bajo desbalance de
clases}
\label{sec:brdl}

Publicado como preprint en arXiv~\cite{su2025brdl}, aborda la pregunta que plantea
\textbf{OE2}: cómo se ajusta un sistema que debe incorporar clases sin supervisión
completa sin degradar las que ya clasificaba bien. Combina dos estrategias: una destilación residual y un aprendizaje balanceado de
pseudo-etiquetas. BRDL trabaja en aprendizaje incremental, donde las clases nuevas sí llegan
etiquetadas cuando les toca su turno, y en el descubrimiento de clases novedosas
nunca lo están.

Sus resultados interesan en tres bloques. Primero, la grilla conjunta de
$\lambda_2$ y $\lambda_3$: el máximo es 46,58 en $(0{,}2;\,0{,}5)$ y el mínimo,
45,1, de modo que toda la grilla cabe en 1,48 puntos de mIoU
\cite[Tabla~4]{su2025brdl}. Segundo, la ablación de las dos estrategias
(Tabla~\ref{tab:ablacion-brdl}). Tercero, el entrenamiento multi-paso, donde una clase se queda en $0{,}00$ de mIoU
en los cinco pasos y para los dos métodos comparados \cite[Tabla~5]{su2025brdl}.

\begin{table}[!htbp]
\centering
\setstretch{1}\small
\caption[Ablación de BRDL en S3DIS]{Ablación de las dos estrategias de BRDL sobre S3DIS, partición S0,
según la Tabla~6 de \cite{su2025brdl}.}
\label{tab:ablacion-brdl}
\begin{tabularx}{\linewidth}{@{}>{\raggedright\arraybackslash\hsize=1.00\hsize}Xcc@{}}
\hline
Configuración & mIoU (\emph{all}) & Diferencia \\
\hline
Sin ninguna de las dos estrategias & 45,19 & referencia \\
Solo una de las dos estrategias & 44,90 & $-0{,}29$ \\
Solo la otra estrategia & 46,83 & $+1{,}64$ \\
Las dos estrategias juntas & 47,26 & $+2{,}07$ \\
\hline
\end{tabularx}
\end{table}

Sirve por tres razones. La primera es metodológica: es el único trabajo revisado
que publica una grilla conjunta de dos parámetros acoplados y la lee como tal.
Concluye que $\lambda_2$ domina y que $\lambda_3$ solo introduce ajustes finos
según el valor de $\lambda_2$. La Tabla~\ref{tab:ablacion-brdl} muestra lo mismo
desde los componentes: activar una sola de las dos piezas empeora el resultado
respecto de no activar ninguna, y solo la combinación gana.

La segunda es de magnitud: la grilla mueve poco frente a los 2,07 puntos que
mueve activar o desactivar los componentes. La calibración es necesaria pero no es donde
está la ganancia grande, y eso fija la expectativa realista de OE2.

La tercera tiene que ver con el desbalance. Con soporte insuficiente el resultado
no mejora de forma gradual: la clase que queda en cero no la recupera ningún valor
de hiperparámetro. Eso coincide con el régimen que esta tesis declara
desfavorable, el de clases novedosas por debajo del 0,5\,\% de los puntos de la
sección. Por esa diferencia de escenario, de BRDL se traslada el procedimiento y
el diagnóstico, no las cifras.

% =====================================================================
\section[Conclusión: aportes del proyecto]{Conclusión de la revisión crítica}
\label{sec:aporte}

Cada estudio previo aborda \textbf{una parte} del problema, pero ninguno lo cubre
por completo:

\begin{itemize}\setlength{\itemsep}{2pt}
\item \textbf{HyperNCD} descubre clases novedosas con un hipergrafo, pero su
descriptor geométrico es de fórmula fija y aporta menos de la mitad que el
hipergrafo (Sección~\ref{sec:hyperncd}).
\item \textbf{Sonata} ofrece una representación de calidad demostrada, pero se
midió en régimen supervisado y con un vocabulario cerrado
(Sección~\ref{sec:sonata}).
\item \textbf{PTv3} muestra cómo opera esa representación: sobre una secuencia
unidimensional, distinta de la rejilla que usa HyperNCD
(Sección~\ref{sec:ptv3}).
\item \textbf{BRDL} muestra cómo calibrar mecanismos acoplados, pero con
parámetros de un mismo método y de una misma escala (Sección~\ref{sec:brdl}).
\end{itemize}

Además, ninguno se evaluó sobre patrimonio construido, y esta revisión no
encontró un trabajo que combine estas piezas. Sobre esa base el proyecto declara
\textbf{siete aportes}, agrupados por objetivo específico
(Tabla~\ref{tab:obj-vs-revision}); el Capítulo~\ref{cap:resultados} y los
trabajos futuros indican cuáles se ejecutaron.

\medskip
\noindent\textbf{OE1. Integración de Sonata en el módulo de prototipos.}
Tres aportes: sustituir el descriptor geométrico (A1), corregir la base de
comparación (A2) y resolver el puente entre las dos estructuras espaciales (A4).

\begin{enumerate}
\item[\textbf{A1}] \textbf{Sustitución del descriptor geométrico.} Reemplazar
$f_{\text{geo}}$ por $f_{\mathrm{ssl}}$, el codificador congelado de Sonata,
conservando íntegro el razonamiento por hipergrafo. Así,
$F_i = [f_{\mathrm{sem}};\, f_{\mathrm{ssl}}]$ y cualquier variación de
la métrica es atribuible a ese reemplazo y a nada más.

\item[\textbf{A2}] \textbf{Auditoría de la implementación de referencia y
corrección de la línea base.} La inspección del código liberado muestra
que el módulo que reproduce la formulación publicada del razonamiento por
hipergrafo no se invoca desde ninguna ruta de entrenamiento ni de evaluación; que
sus matrices de peso se crean dentro de la función en cada llamada, por lo que no
reciben gradiente ni se registran como parámetros; que la llamada actualización
dinámica repite la construcción inicial sobre la misma matriz de similitud y no
introduce ninguna actualización; y que el descriptor geométrico ejecutado tiene
dimensión seis, por el centroide de la región, y no tres como describe el
artículo. De ahí tres decisiones: la línea base se establece sobre el código que
se ejecuta y no sobre la formulación publicada, la comparación se hace contra
corridas propias en condiciones idénticas, y la sustitución se aplica sobre el
descriptor de dimensión seis. Sin esta corrección, la mejora no se podría atribuir
con certeza a la sustitución.

\item[\textbf{A4}] \textbf{Puente entre dos estructuras espaciales
incompatibles.} Resolver el emparejamiento entre la secuencia serializada de PTv3
y la rejilla dispersa de Minkowski mediante \emph{up-cast} de resolución y
\emph{cache} de características por sección. La misma decisión resuelve tres
problemas: el emparejamiento de vecindades, la incompatibilidad entre pilas de
dependencias (deja de importar porque las dos ramas nunca coexisten en un
proceso) y la repetición del costo del codificador en cada época. Incluye una cota
de aceleración propia que delimita hasta dónde vale la pena distribuir el
cacheo.
\end{enumerate}

\medskip
\noindent\textbf{OE2. Calibración del híbrido.} Tres aportes: normalizar la
escala entre ramas (A5), calibrar de forma conjunta los parámetros de la afinidad
(A6) y verificar la configuración fuera de la sección de calibración (A7).

\begin{enumerate}
\item[\textbf{A5}] \textbf{Normalización de escala entre dos ramas
heterogéneas.} El descriptor híbrido $F_i=[f_{\mathrm{sem}};f_{\mathrm{ssl}}]$ une dos
ramas de escala distinta, con 96 componentes entrenadas junto con el resto del
modelo y 1088 componentes congeladas cuyo rango no está normalizado contra el de
la rama semántica. Sin corrección, la rama más ancha domina la similitud por
coseno $S_s$ con la que se comparan las regiones, y el valor publicado
$\alpha=0{,}5$ deja de repartir el peso entre forma y significado como en la
línea base. Hacer comparables las dos ramas antes de concatenarlas es un
requisito previo a la calibración, y el régimen congelado lo agrava porque impide que la rama auto-supervisada se reescale sola
durante el entrenamiento.

\item[\textbf{A6}] \textbf{Calibración conjunta y verificación de
no-aditividad.} Determinar los tres valores con una grilla
conjunta y curvas de sensibilidad sobre una sección de referencia, en lugar de un
barrido por parámetro, y verificar si la no-aditividad que BRDL documenta entre
sus dos estrategias aparece también entre la representación y el razonamiento del
híbrido.

\item[\textbf{A7}] \textbf{Verificación de que la configuración se sostiene fuera
de donde se calibró.} La configuración se fija sobre una única sección de
referencia y se aplica \textbf{sin reajuste} a las restantes, reportando
la degradación observada. Así se distingue una calibración real de un ajuste hecho
sobre aquello que se iba a medir.
\end{enumerate}

\medskip
\noindent\textbf{OE3. Cuantificación de la mejora y delimitación del rango de
operación.} Un aporte: formular una métrica que separe la contribución de la
representación de la del razonamiento, con su protocolo de atribución por etapas
(A3).

\begin{enumerate}
\item[\textbf{A3}] \textbf{Métrica de brecha cerrada $G$ y protocolo de
atribución por etapas.} Las métricas estándar no separan lo que aporta la
representación de lo que aporta el razonamiento: el mIoU final es el mismo número
venga de donde venga la ganancia. Trasladando el principio del \emph{linear
probing}, este trabajo formula la brecha cerrada $G$ (Ec.~\eqref{eq:G}), que
expresa qué parte de la distancia entre clases conocidas y novedosas se logró
recuperar, y mide también en el punto intermedio del proceso
(Sección~\ref{sec:metricas-hibrido}). Incorpora además un análisis de dominancia
que identifica la clase que gobierna cada mejora agregada.
\end{enumerate}


